\section{Results}

The replicated GSE168041 profiles comprised 64,705 cells from 16 samples, including 31,556 spiral ganglion and 33,149 lateral-wall cells (Fig.~\ref{fig:1}; Supplementary Fig. S1). Thirteen tissue--cell-type combinations retained four unexposed and four noise-exposed samples. Fibrocyte/\allowbreak{}mesenchymal, endothelial and immune/\allowbreak{}myeloid populations were represented in both compartments. Sample-level retention after quality control and doublet removal ranged from 86.6\% to 92.6\%; differential-expression counts for all thirteen profiles are shown in Supplementary Fig. S2.

\subsection{NF-\texorpdfstring{\(\kappa\)}{kappa}B responses diverge across cochlear cell types}

TNF--NF-\(\kappa\)B enrichment differed in direction across cochlear cell types (Fig.~\ref{fig:2}). Fibrocyte/\allowbreak{}mesenchymal cells showed enrichment toward noise exposure in the lateral wall (NES 1.648; FDR 0.00371) and spiral ganglion region (NES 1.463; FDR 0.0305). In the lateral wall, the 70-gene leading edge included the chemokine gene \mbox{\textit{Ccl5}}, which had a \(\log_2\) fold change of 5.958 and FDR 4.12 \(\times\) \(10^{-14}\). Other leading-edge members included \mbox{\textit{Ifih1}}, \mbox{\textit{Mcl1}} and \mbox{\textit{Ifit2}}, with \(\log_2\) fold changes of 1.321, 0.467 and 2.061, respectively, and FDR values of 0.00227, 0.0209 and 0.0375.

Endothelial and immune/\allowbreak{}myeloid populations showed enrichment of the same gene set toward the unexposed condition in both regions. In the lateral wall, endothelial cells had NES \(-\)2.813 and FDR 4.27 \(\times\) \(10^{-20}\), and immune/\allowbreak{}myeloid cells had NES \(-\)3.213 and FDR 1.78 \(\times\) \(10^{-32}\). In the spiral ganglion preparation, the corresponding NES values were \(-\)1.854 and \(-\)2.375, with FDR values of 7.62 \(\times\) \(10^{-5}\) and 1.38 \(\times\) \(10^{-12}\). Positive TNF--NF-\(\kappa\)B enrichment also occurred in spiral ganglion neurons (NES 1.820; FDR 5.17 \(\times\) \(10^{-5}\)) and lateral-wall epithelial/\allowbreak{}supporting-associated cells (NES 2.143; FDR 1.31 \(\times\) \(10^{-8}\)).

Lateral-wall IL6--JAK--STAT3 signaling was enriched toward the unexposed condition in endothelial cells (NES \(-\)1.796; FDR 0.00331) and immune/\allowbreak{}myeloid cells (NES \(-\)1.952; FDR 0.000362). Endothelial MAPK family signaling cascades were enriched toward noise exposure (NES 1.918; FDR 8.37 \(\times\) \(10^{-7}\)). Its 48-gene leading edge contained 28 proteasome-subunit genes, together with \mbox{\textit{Map2k2}}, \mbox{\textit{Mapk3}} and \mbox{\textit{Dusp2}}. Endothelial \mbox{\textit{Map2k2}} and \mbox{\textit{Mapk3}} had \(\log_2\) fold changes of 0.417 and 0.336, with FDR values of 0.000427 and 0.00466. Endothelial angiogenesis and VEGF enrichment did not meet the FDR criterion, with NES values of 1.173 and \(-\)1.061 and FDR values of 0.384 and 0.513.

\subsection{Opposing enrichment directions persisted after individual-animal omission}

All 56 focused pathway estimates retained the direction of the full-sample analysis across the 48 leave-one-animal-out fits (Fig.~\ref{fig:3}). Fibrocyte/\allowbreak{}mesenchymal TNF--NF-\(\kappa\)B enrichment remained positive in all eight omissions in each region. Spiral-ganglion estimates ranged from NES 1.305 to 1.731, with FDR \(<\) 0.05 in seven of eight fits. Lateral-wall estimates ranged from NES 1.317 to 2.007, with six of eight fits meeting this criterion. FDR ranges were 0.000326--0.136 and 3.94 \(\times\) \(10^{-6}\)--0.169, respectively. Endothelial and immune/\allowbreak{}myeloid TNF--NF-\(\kappa\)B enrichment remained negative and significant in all eight omissions in each region.

Lateral-wall endothelial MAPK-family enrichment remained positive and significant in all eight omissions, with NES 1.807--1.925 and FDR 5.50 \(\times\) \(10^{-7}\)--2.08 \(\times\) \(10^{-5}\). Individual-animal pathway profiles are shown alongside the omission estimates in Fig.~\ref{fig:3}; expression profiles for the reported leading-edge genes are shown in Supplementary Fig. S3.

\subsection{Whole cochleae showed inflammatory and proteostatic enrichment at both exposure levels}

Whole-cochlea profiles comprised 19,673 genes from twelve samples, with four samples at each of 70, 94 and 105 dB (Fig.~\ref{fig:4}). At FDR \(<\) 0.05, 5,777 genes differed between 94 and 70 dB, 889 between 105 and 70 dB, and 4,177 between 105 and 94 dB (Supplementary Fig. S4). The overall dose likelihood-ratio test identified 6,398 genes.

TNF--NF-\(\kappa\)B, protein-folding and unfolded-protein-response gene sets were positively enriched at both 94 and 105 dB relative to 70 dB (Fig.~\ref{fig:4}). Overall, 17 of 27 pathway--contrast combinations met global FDR \(<\) 0.05: eight for 94 versus 70 dB and nine for 105 versus 70 dB. For TNF--NF-\(\kappa\)B signaling, the corresponding NES values were 2.563 and 2.979, with global FDR values of 3.45 \(\times\) \(10^{-18}\) and 4.41 \(\times\) \(10^{-24}\). At 94 dB, the unfolded protein response had NES 2.059 and global FDR 3.75 \(\times\) \(10^{-6}\). At 105 dB, this pathway had NES 2.512 and global FDR 5.17 \(\times\) \(10^{-10}\). Protein folding yielded NES 1.800 and 2.155, respectively, with global FDR values of 2.80 \(\times\) \(10^{-5}\) and 3.06 \(\times\) \(10^{-8}\).

None of the nine targeted pathways met the global FDR criterion in the direct 105-versus-94-dB comparison. TNF--NF-\(\kappa\)B signaling yielded NES \(-\)1.135 and global FDR 0.223, while the reactive oxygen species pathway yielded NES \(-\)1.488 and FDR 0.0595. The two exposed groups shared positive inflammatory enrichment relative to the 70-dB reference, while retaining distinct genome-wide profiles. Positive whole-cochlea TNF--NF-\(\kappa\)B enrichment accompanied the fibrocyte-associated response observed in the separately sampled GSE168041 dataset.

\subsection{Eleven selected genes retained endothelial or fibrocyte localization across three libraries}

Eleven of the twenty selected genes had the same highest-expression broad cell category across all three GSE252285 libraries (Fig.~\ref{fig:5}; Supplementary Fig. S5). Categories were ranked by mean normalized expression, and each corresponding group contained at least 20 cells. The 36,111 cells comprised 10,273 control, 16,631 hidden hearing loss and 9,207 noise-induced hearing loss profiles.

The seven genes with recurrent endothelial expression maxima were \mbox{\textit{Dusp1}}, \mbox{\textit{Nfkbia}}, \mbox{\textit{Apold1}}, \mbox{\textit{Sik1}}, \mbox{\textit{Fosb}}, \mbox{\textit{Rasd1}} and \mbox{\textit{Depp1}}. The four genes with recurrent fibrocyte/\allowbreak{}mesenchymal expression maxima were \mbox{\textit{Ccn1}}, \mbox{\textit{Errfi1}}, \mbox{\textit{Nr4a1}} and \mbox{\textit{Egr1}}. Endothelial \mbox{\textit{Nfkbia}} detection fractions were 87.5\%, 85.3\% and 87.3\% in control, hidden hearing loss and noise-induced hearing loss libraries, respectively. Fibrocyte/\allowbreak{}mesenchymal \mbox{\textit{Egr1}} detection fractions were 89.5\%, 91.5\% and 87.3\%.

The eleven recurrent cell-type localizations persisted when cluster 19 was included in the fibrocyte/\allowbreak{}mesenchymal group (Fig.~\ref{fig:6}; Supplementary Fig. S6). \mbox{\textit{Cebpb}} additionally acquired a recurrent fibrocyte/\allowbreak{}mesenchymal maximum. Cluster 19 expressed fibrocyte-associated markers, including \mbox{\textit{Coch}} in 99.7\% and \mbox{\textit{Otor}} in 82.1\% of its 319 cells. \mbox{\textit{Pecam1}} and \mbox{\textit{Kdr}} were detected in 1.9\% and 9.1\%, respectively, and \mbox{\textit{Vwf}} was not detected.

\mbox{\textit{Apold1}} was highly expressed in both endothelial cells and cluster 19 (Supplementary Fig. S7). Its mean normalized expression was 1.585 in cluster 19 and 1.338 across endothelial cells, with detection fractions of 75.9\% and 66.9\%, respectively. Endothelial cells retained the highest broad-cell-category mean in all three libraries after cluster 19 was included in the fibrocyte/\allowbreak{}mesenchymal group.

Five of the six UCell-guided follow-up comparisons exceeded the expression- and detection-matched null at the six-test FDR threshold (Supplementary Fig. S8). The endothelial and fibrocyte/\allowbreak{}mesenchymal TNF--NF-\(\kappa\)B localization statistics exceeded the matched null distributions, with observed score differences of 0.0545 and 0.0380 and FDR 0.000200 for both tests. Endothelial VEGF/\allowbreak{}angiogenesis and endothelial and immune/\allowbreak{}myeloid IL6--JAK--STAT3 localization also met the six-test FDR criterion. The endothelial MAPK/\allowbreak{}AP-1 localization statistic yielded a score difference of 0.0176 and FDR 0.0724. Adding gene-length strata retained the five primary-positive localization results. Reassigning cluster 19 preserved all six significance classifications in both matching analyses (Supplementary Fig. S6). The fibrocyte/\allowbreak{}mesenchymal TNF--NF-\(\kappa\)B score difference increased from 0.0380 to 0.0394, with FDR 0.000200 in both annotations. Endothelial MAPK/\allowbreak{}AP-1 remained nonsignificant, with primary FDR 0.0725 and gene-length-matched FDR 0.0817 under the alternative annotation.

In GSE168041, lateral-wall endothelial \mbox{\textit{Apold1}} had a \(\log_2\) fold change of 0.031 and FDR 0.913. Fibrocyte/\allowbreak{}mesenchymal \mbox{\textit{Errfi1}} had a \(\log_2\) fold change of \(-\)0.361 and FDR 0.588.
